\documentclass[11pt]{article}
\usepackage[margin=1in]{geometry}
\usepackage{amsmath,amssymb,hyperref}
\newcommand{\F}{\mathbb F}
\newcommand{\ord}{\operatorname{ord}}
\title{The absolute unit and two exact moments at one target}
\author{Project research note}
\date{18 November 2026}
\begin{document}
\maketitle

On the project's odd prime-order family one can compute the local
unit $U_t$ and both integer moments from the dyadic seeds at any
nonidentity target. A balanced split uses $2^{(1/2+o(1))b}$ work
and space, including seed expansion. This covers arbitrary target
multiplicity. A second construction reduces the same data to counts
of individual bits with a fixed endpoint and no table of either
half. Its inexpensive specialization recovers low-order
multiplicative characters of the unit and the exact integer moments.
The full-field specialization has a more expensive arithmetic
circuit. Neither construction improves the project's conditional
end-to-end exponent.

\section{Local elimination with multiplicities}
Take $E:y^2=x^3+ax+c$ over $\F_p$, $p>3$,
$T=(t,y_T)$ with $y_T\ne0$, and paired affine abscissa sets
$U_i$ of sizes $n_i$. Put $b=\lceil\log_2 p\rceil$.
For window $(\alpha_i,\beta_i,\tau_i)$ modulo $M=2^b$, enumerate
$x=\alpha_i^{-1}(j-\beta_i)\bmod M$ for $0\le j<\tau_i$,
require $x<p$, and take one square root of the nonzero curve
right-hand side if one exists. This costs
$\sum_i\tau_i\,\mathrm{poly}(b)$ expected bit operations with
ordinary finite-field square roots. There is no group-index table.
An empty slot has order zero, unit one and zero moments.

Write $N(z)$ for the normalized product of summation forms in
\texttt{../scalar\_norm/norm.tex}. The requested unit is the
coefficient of $(z-t)^h$, for $h=\ord_t N$. The moments are sums of
$q$ and $q^2$ over signed tuples summing to $T$, with the canonical
integer $q=\sum_{i=1}^m u_i p^{i-1}$. This agrees with folding
the two target signs in the oriented convention of the section note.

Split the slots into a nonempty left block and a right block.
The left records run over abscissas and signs modulo simultaneous
negation, with first sign fixed. Store $(P,q_L)$ for each left
sum. The right records are $(Q,q_R)$ with $Q=T+\sum_R\epsilon_iP_i$
and all right signs free. The two codes use the original digit
positions. Their respective list lengths are
\[
 f=2^{|L|-1}\prod_{i\in L}n_i,\qquad
 g=2^{|R|}\prod_{i\in R}n_i .
\]
Use the chord scale $L(P,V)$ of \texttt{../section\_product/sections.tex}.
For each appended slot of size $n$, update the common scalar of
a half by $\beta\gets\beta^{2n}\prod_{P,V}L(P,V)$ while extending
each root to $P+V,P-V$. A finite root contributes $W_0-x(P)W_1$,
and an infinite root contributes $-W_1$.
Let $k_u,l_u,s_u$ be the count and two code sums
$(\sum q_L,\sum q_L^2)$ in each left abscissa bucket, including
one for infinity. The affine monic polynomial and leading scalar are
\[
 F(W)=\prod_{u\text{ finite}}(W-u)^{k_u},\qquad
 a_L=(-1)^{k_\infty}\beta_L.
\]

At a distinct finite right abscissa $u$ compute
$F\bmod(W-u)^{\max(2,k_u+1)}$ by a product and remainder tree.
If $k_u>0$, its coefficient of $W^{k_u}$ is
$v_u=(F/(W-u)^{k_u})(u)$. Otherwise it gives both
$v_u=F(u)$ and $F'(u)$. This step performs no division by a
factor that vanishes at $t$. The total modulus degree is at most
$2g+f$, irrespective of the target order.

For every finite right record use $e_Q=k_{x(Q)}$ and
$u_Q=v_{x(Q)}(y(Q)/y_T)^{e_Q}$. For an infinite right record
use $e_Q=2k_\infty$ and $u_Q=(4y_T^2)^{-k_\infty}$. Then
\[
 h=\sum_Q e_Q,\qquad
 U_t=\beta_R^f a_L^g\prod_Q u_Q .
\]
To prove the formula, deform the target abscissa $z$ at $t$.
The normalized split resultant is the product of the left binary
form evaluated on the moving right homogeneous root vectors, times
$\beta_R(z)^f$. The interchange sign is positive because $g$ is
even for a nonempty right block. With one original slot use
$N(z)=\prod(z-u)$ directly, which gives the same displayed formula
with the sole right root $T$. At a finite specialized right root
use coordinates $(x(Q(z)),1)$. The invariant differential
$dx/(2y)$ gives $dx(Q(z))/dz=y(Q)/y_T$. At a right root
specializing to infinity use $(1,1/x(Q(z)))$. In the parameter
$s=-x/y$ at infinity, $1/x=s^2+O(s^3)$ and
$dx/(2y)=(1+O(s))ds$. Consequently the reciprocal coordinate
starts with $(z-t)^2/(4y_T^2)$. These are exactly the leading
coefficients in the formula. The differential and formal parameter
are described in III.5 and IV.1 of
\href{https://www.math.brown.edu/johsilve/AECHome.html}
{Silverman's \emph{The Arithmetic of Elliptic Curves}}.

For the integer moments, for each right record of code $q_R$
look up its left bucket $(k,l,s)$. Its contributions are
\[
 \mu(k,\ l+kq_R,\ s+2q_Rl+kq_R^2),
 \qquad \mu=1\ \text{at a finite root},\quad \mu=2\ \text{at infinity}.
\]
Indeed a signed solution of total $T$ uses the right signs negated
and precisely one of the two left orientations for a finite match.
Both left orientations work at infinity. Thus the same join gives
$h,M_1,M_2$ even if many records have identical abscissas. At a
zero use $hM_2-M_1^2$ over the integers as in the preceding note.
Positive variance makes all first derivatives zero. Zero variance
decodes the single abscissa tuple from $M_1/h$, and dividing
$U_t$ by that tuple's leading local coefficient gives its full
additive cofactor.

At order one $N_z=U_t$, and at higher order $N_z=0$. At order zero
differentiate the nonvanishing product. The logarithmic derivative
is $f\beta_R'/\beta_R+\sum_{Q\ne O}(y(Q)/y_T)F'(x(Q))/F(x(Q))$.
The moving half propagates its scale derivative through each chord.
The logarithmic increment for appending $V$ to a finite moving
root $P$ with distinct abscissa is
$2y(P)/(y_T(x(P)-x(V)))$. At $P=\pm V$ it is
$(3x(P)^2+a)/(2y(P)y_T)$.
Indeed switching the infinite outgoing root to its reciprocal chart
replaces $(x(P)-x(V))^2$ by
$(y(P)+y(\pm V))^2-(x(P)+x(V))(x(P)-x(V))^2$, taking
$y(\pm V)=y(P)$ at specialization. An incoming infinity gives
zero increment in its reciprocal chart. These formulas account
for chart changes even on a nonvanishing target. The literal
additive marker derivatives at a nonzero still have the separate
two-slot algorithm of the scalar note.

Record generation and the single remainder tree cost
$(f+g)\,\mathrm{poly}(b,m,\log(f+g))$ expected bit operations up
to logarithmic factors, in addition to slot extraction and the
optional $2^{O(m)}\mathrm{poly}(b,m)$ isolated-tuple form.
For the tuned widths and a balanced split this is
$2^{(1/2+o(1))b}$ including storage up to polynomial factors.
This extends the previous half-split calculation to the actual
higher-order unit and integer moments. Its explicit lists have
$fg=2^{m-1}\prod n_i$. On the tuned family this is of order $r$ in
probability by the earlier slot-size estimate, where $r=\#E(\F_p)$.

\section{A joint counting reduction with no half lists}
Define $w_T(S)=t-x(S)$ at finite points away from $\pm T$,
$w_T(\pm T)=1$ and $w_T(O)=-1$. In the forward traversal fix
the first sign, write $\mathcal P_j$ for the multiset of oriented
prefix sums, and take one lift $V$ per abscissa in a new slot.
The same normalized chord recurrence proves
\[
 U_t=\prod_{S\in\mathcal P_m}w_T(S)
       \prod_{j=2}^m
       \left(\prod_{P\in\mathcal P_{j-1},\,V\in U_j}
              L(P,V)\right)^{\prod_{k>j}2n_k}.
\]
Here $V\in U_j$ denotes the chosen lift. The last finite linear
factors give the first product after removing exactly the target
zeros. Each earlier scale is raised to $2n_k$ on continuing to
slot $k$. This expression avoids evaluating the function
$\gamma$ on any point. The products index mathematical
contributions. The following counting reduction does not generate
the prefix lists.

For $d\mid p-1$ choose $g$ whose image generates the quotient by
$d$th powers, and set $\zeta=g^{(p-1)/d}$. Every nonzero local
weight $w$ has one class $0\le e<d$ for which
\[
                     g^e v^d=w,\qquad v\in\F_p^\times .
\]
There are exactly $d$ witnesses $v$. Pin each bit of the canonical
$e$ to one in turn and count its occurrences. Dividing the
counts by $d$ and combining them with binary place values gives
the sum of the class labels. Work with all signs. Divide the
terminal sums by two and the prefix sums by four in the integers
before reducing modulo $d$. The global flip and the independent
sign of the new slot preserve the respective weights. Include the
prefix exponents from the displayed product. If the resulting sum
is $s$ modulo $d$, the answer is $U_t^{(p-1)/d}=\zeta^s$.
For $d=2$ every chord scale is a square, so only the terminal
predicate is needed. For $d=p-1$, omit $v$ and its multiplicity.
The certificate becomes $g^e=w$ and yields the entire $U_t$.
The bits require only $O(m\log d)$ kinds of count. Enumerating all
$d$ classes gives another option for tiny $d$.

The moments use only the ordinary count at $T$ and bit pins on the
existing canonical abscissas. Write $q=\sum_\ell a_\ell b_\ell$
where $a_\ell=p^{i-1}2^k$ and $b_\ell$ is bit $k$ of $u_i$.
Then $M_1=\sum a_\ell H_\ell$ and
$M_2=\sum a_\ell^2H_\ell+
2\sum_{\ell<s}a_\ell a_s H_{\ell s}$. Single and double pins
give these counts in $O((mb)^2)$ calls. No rounding is applied
to the large integer moments themselves.

\paragraph{Implementing the predicates in the fixed-endpoint graph.}
Use the offset construction in \texttt{../short\_hash\_count/count.tex}.
For a prefix (or all slots) of signed sum $S\ne O$, extend to
the guarded path
\[
                D,A_1,\ldots,A_j,B,H,-D\quad\hbox{ending at }H
\]
for a public nonidentity $H$. Cancellation forces the unique
affine $B=-S$. The terminal predicate sees $x(B)$, including
the patched target value. The analogous prefix predicate also
checks the next independent signed point $V$ with its ordinary
curve equation and hash. Its weight is $(x(B)-x(V))^2$ for
distinct abscissas and $4y(V)^2$ on equality. The identity
contribution is handled by the existing fixed-identity count. Its
terminal weight is $-1$ and its prefix weight is one. Fixed
endpoint membership keeps the entire guarded path on the supplied
curve. Comparisons, the independent point and the constant extra
slots cost $O(b)$ long-word edges in the existing compiler.

For each unguarded assignment at most $c_m=2m+11$ choices of the
offset are inadmissible. The initial identity costs one, the
interior guards exclude at most two offsets per step, and the
last guard excludes at most four because $H\ne O$. The auxiliary
root $v$ is independent of the guards. If $Z$ is a raw count for
a bit pin, divide out its fixed root multiplicity first. The
quotient $Z'$ satisfies $0\le rC-Z'\le c_m C$ for the desired
signed count $C$. A class predicate can cover a positive fraction
of all endpoints, so rounding on the original box need not work.
Split dyadic windows until $c_m\prod_i(2\tau_i)<r$ on each box.
Then $C=\lceil Z'/r\rceil$ separately on every box. At the tuned
and paid widths $\sum_i\log_2(2\tau_i)\le b+O(m\log b)$.
Binary splitting therefore uses at most $2^{O(m\log b)}$ boxes.
There are enough free index bits since $c_m2^m<r$ eventually.
All split offsets and bit pins are public, and every surviving
tuple is counted in exactly one box. The finite guard test is
part of the reduction.

\paragraph{Cost of obtaining the full residue.}
Binary powering with a bounded exponent witness and a fixed
exponent on $v$ uses $O(\log d)$ modular product equations, each
with $O(b)$ word edges. For the public base, the constants
$g^{2^i}$ can be prepared ahead of time. Each selection and
intermediate residue is forced by the exponent bits. For a small
$d$, trial division of $d$ and a random base suffice to check that
$\zeta$ has order $d$, with no factorization of $p-1$. For the
full bit reduction a primitive base can be supplied with a
verified order certificate, or selected by the same full-width
counting mechanism. Exactly one $0\le e<p-1$ solves $g^e=1$
precisely for a primitive base. Uniform sampling takes $O(b)$
trials in expectation, since $\varphi(n)/n\ge1/(1+\log_2 n)$
by ordering the distinct prime divisors.

For $\log_2 d=o(m)$ the additions cost $o(mb)$ word edges, with
at most polylogarithmically many new arrays in the tuned regime.
Use canonical range, quotient, carry and local bit checks of the
original compiler. For a variable exponent bit, a canonical
auxiliary field input is either one or the appropriate public
power of $g$. Copy the bit to the existing source-row owners and
check their owned limbs locally. This costs only
$\sqrt b\,\mathrm{polylog}(b)$ auxiliary binary edges per bit.
The long arithmetic words still share endpoints for alphabet
recoding, and the short carries retain the earlier lower-order cost.
The original degree and parallel-multiplicity bounds are preserved
by the same row grouping and relays. The polynomially many pins
and the logarithm of the number of boxes are lower order.
Consequently the conditional leading query estimate is still
$(18K+o(1))mb/\sqrt{\log_2 b}$ in log time, with the $K$ and
finite-packet assumptions of the short-hash note. This establishes
an inexpensive joint reduction for the moments and those
characters. The character-bit rank obstruction of the section note
applies to separated formulas and permits this joint count.

For $d=p-1$ the displayed binary certificate has $\Theta(b)$
stages and the added long-word edge budget is $\Theta(b^2)$ in
this implementation. The earlier leading query estimate does not
apply. Using several cheap orders gives the projection of order
their least common multiple. If all orders are at most $D$, their
least common multiple is at most $D!$. Covering $p-1$ therefore
requires $D\log_2 D\ge b-O(1)$ and $\log_2 D=\Omega(\log b)$.
At tuned $m=\Theta((\log b)^{1/4})$ this is outside the
no-leading-overhead regime of this binary circuit. This accounting
does not give a lower bound for arbitrary arithmetic circuits,
adaptive queries, or algorithms operating directly on the seeds.
The best end-to-end estimate in the project remains conditional on
the original finite-packet and contraction assumptions,
\[
 \log_2\mathbb E T_{\rm solve}\le(26.2933775156+o(1))
                       \frac{b}{(\log_2 b)^{1/4}} .
\]

\section{Reproducible checks}
Run \texttt{python -B research/single\_target/verify.py} and
\texttt{python -B research/single\_target/covers\_check.py} with
python-flint. The finite verifier uses word-size prime fields with
\texttt{nmod\_poly}. The complexity estimates use general finite-field
polynomial arithmetic. The first check compares 603 small inputs and targets with
literal tuple products and signed integer moments, including
maximum order 35 in characteristic 13. All 117 zero gradients
also matched. On the 2,800-tuple fixture the two lists have 56
and 400 entries. At $(0,1)$ the order is one, the unit is 43446,
and the first moment is 9199523451172809648. The second is its
integer square. The second checker exhausts guarded offsets on
the order-223 curve over $\F_{211}$. With $210=2\cdot3\cdot5
\cdot7$, the four projections and the direct exponent-bit counts
independently reconstruct the unit, including a target of order
21. The finite checker exercises the count identities, while
the asymptotic count cost uses the project's contractor assumption.

At $T=(72,189)$ the boxes $((1,0,2),(1,205,8))$ and
$((1,0,2),(11,185,8))$ give the same slot sizes $(1,5)$, order
one, $M_1=11183$ and $M_2=125059489$. Both order-30 projections
are 156. Their full units are 104 and 200. Even on an isolated hit
the missing component changes the absolute derivative.

\paragraph{Follow-up on full-unit recovery.}
The companion note \texttt{unit\_recovery.tex} gives a direct divisor-count
method and a uniform block certificate with $O(b^{4/3})$ additional
Boolean edges. It also bounds the number of exceptional values the
divisor method must visit on large occupied slots. The conditional
end-to-end coefficient remains $26.2933775156$.
\end{document}
